\documentclass{article}

\usepackage{tabularx}
\usepackage{booktabs}

\title{Problem Statement and Goals\\\progname}

\author{\authname}

\date{}

\input{../Comments.text}
\input{../Common.text}

\begin{document}

\maketitle

\begin{table}[hp]
\caption{Revision History} \label{TblRevisionHistory}
\begin{tabularx}{\textwidth}{llX}
\toprule
\textbf{Date} & \textbf{Developer(s)} & \textbf{Change}\\
\midrule
September 24th, 2026 & Pedro & Initial draft of document from discussion\\
Date2 & Name(s) & Description of changes\\
... & ... & ...\\
\bottomrule
\end{tabularx}
\end{table}

\section{Problem Statement}

\subsection{Problem}

The McMaster Troitsky club is a student-run organization that registers and 
supports 4-6 teams of students as they take on the task of modelling, analyzing 
and manfacturing a bridge out of popsicle sticks, white glue and dental floss. 
While the club has been successful in the past, results often ranges in performance, 
both between years and teams, due to the lack of centralized information and tehcnical 
support for members. As teams progress through several years of competiton, 
undergoing material testing, bridge design and analysios, and even experimenting with 
more innovative concepts, skill and knowledge within the club often consolidates within 
more experienced members or single teams, ensuring that when those members graduate information 
is lost to the club. This project seeks to create a centralized information repository and
technical support system for the McMaster Troitsky club, ensuring information transfer between years and teams, 
and facilitating support and onboarding for new members.

\subsection{Inputs and Outputs}

\textbf{Inputs:} WHAT WOULD THE INPUTS BE? IS IT ABOUT THE INPUTS INTO THE PROJECT OR THE INPUTS INTO THE SYSTEM?. 
Range of McMaster Troitsky club member and alumi knowledge and past documentation of club and team specific activities.
\\
\textbf{Outputs:} Toolbox of services to support the McMaster Troitsky club in information centralization and technical support for members.

\subsection{Stakeholders}

\subsubsection{Primary Stakeholders}

\begin{itemize}
    \item \textbf{McMaster Troitsky Club General Members}
    \\
    The general members of the club are primary stakeholders as they will be the main 
users and those that we will focus on supporting with our project. These users will be both supporting and expanding the data the
system provides, and will be the main beneficiaries of the system.
    \item \textbf{McMaster Troitsky Club Executive Team}
    \\
    The McMaster Troitsky Club Executive Team are primary stakeholders as they will benefit from 
both the organizational support the service provides, and the aliviation of the burden of technical 
support and information transfer that is currently placed on them. Their input will be highly important 
for service design and prioritization of features.
\end{itemize}

\subsubsection{Secondary Stakeholders}

\begin{itemize}
    \item \textbf{Concordia University's Troitsky Organization}
    \\
    As Concordia University hosts the Troitsky competition, they establish the rules and regulations of the competition, 
and therefor change what is prioritized by the primary stakeholders. This would influence the design of the system, 
and what features would benefit teams the most.
    \item \textbf{MES - McMaster Engineering Society}
    \\
    The MES organizes engineering clubs and seeks to benefit from the legitimization 
of the McMaster Troitsky Club through a centralized information platform. The MES also provides 
funding for the McMaster Troitsky Club, so the hosting of services and costs associated may be dependent on the givenbudget.
\end{itemize}

\subsubsection{Tertiary Stakeholders}

\begin{itemize}
    \item \textbf{McMaster University}
    \\
    The university itself is a tertiary stakeholder as it benefits from the success of the McMaster Troitsky Club, 
and the success of teams during the Troitsky competition.
\end{itemize}

\subsection{Environment}

\subsubsection{Software Environment}
The Troitsky Toolbox will be a web-based application, accessible through any modern web browser. ANYTHING ELSE, ANY HARDWARE ENVIRONMENT?

\section{Goals}

\begin{itemize}
    \item \textbf{Centralize Birdge Development Information} Create a centralized information repository and technical support system for the McMaster Troitsky club.
    \item \textbf{Provide Technical Support} Provide technical support and resources for club members.
    \item \textbf{Facilitate Knowledge Transfer} Facilitate knowledge transfer between years and teams.
\end{itemize}

\section{Stretch Goals}

\begin{itemize}
    \item \textbf{Centralize Birdge Development Information} Create a centralized information repository and technical support system for the McMaster Troitsky club.
    \item \textbf{Provide Technical Support} Provide technical support and resources for club members.
    \item \textbf{Facilitate Knowledge Transfer} Facilitate knowledge transfer between years and teams.
\end{itemize}


\section{Custom Documents (CDs)}

\wss{For CAS 741: State whether the project is a research project. This
designation, with the approval (or request) of the instructor, can be modified
over the course of the term.}

\wss{For SE Capstone: List your custom documents (CDs).  Potential CDs include
usability testing, code walkthroughs, user documentation, formal proof,
GenderMag personas, Design Thinking, etc.  (The full list is on the course
outline and in Lecture 02.) Normally the number of CDs will be two (one per
term).  Approval of the CDs will be part of the discussion with the instructor
for approving the project. The CDs, with the approval (or request) of the
instructor, can be modified over the course of the term.}

\begin{itemize}
    \item \textbf{Interviews}
    \item \textbf{Usability Testing}
\end{itemize}

\newpage{}

\section*{Appendix --- Reflection}

\wss{Not required for CAS 741}

\input{../Reflection.text}

\begin{enumerate}
    \item What went well while writing this deliverable? 
    \item What pain points did you experience during this deliverable, and how
    did you resolve them?
    \item How did you and your team adjust the scope of your goals to ensure
    they are suitable for a Capstone project (not overly ambitious but also of
    appropriate complexity for a senior design project)?
\end{enumerate}  

\end{document}